% arara: pdflatex
\documentclass[12pt]{article}
\usepackage[utf8]{inputenc}
\usepackage[russian]{babel}
\usepackage{wrapfig}
\usepackage{amsmath}
\usepackage{amssymb}
\usepackage{geometry}
\usepackage{graphicx}
\usepackage{rotating}
\usepackage{makecell}
\usepackage{array}
\usepackage{tikzsymbols}
\newcolumntype{P}[1]{>{\centering\arraybackslash}p{#1}}
\newcolumntype{M}[1]{>{\centering\arraybackslash}m{#1}}

\geometry{top=0.5cm}
\geometry{bottom=0.5cm}
\geometry{left=1cm}
\geometry{right=1cm}
\pagestyle{empty}

\usepackage{booktabs}
% заповеди из докупентации:
% 1. Не используйте вертикальные линни
% 2. Не используйте двойные линии
% 3. Единицы измерения - в шапку таблицы
% 4. Не сокращайте .1 вместо 0.1
% 5. Повторяющееся значение повторяйте, а не говорите "то же"

\newcommand{\myrotcell}[1]{\rotcell{\makebox[0pt][l]{#1}}}

\newcommand{\protocol}[6]{
\begin{minipage}{.2\textwidth}
    \begin{center}
    \includegraphics[width=0.48\textwidth]{klshlogo.pdf}
    \end{center}
\end{minipage}
\begin{minipage}{.6\textwidth}
    \begin{flushleft}
    \textbf{Старший судья}: #5\\ 
    \textbf{Судьи}: #6 \\ 
    \textbf{Аудитория}: #2 \\ 
    \end{flushleft}
\end{minipage}
\begin{minipage}{.2\textwidth}
    \begin{flushleft}
        {\Large\textbf{ФМТ тур-#1}} \\
        \textbf{Команды}: #3 --- #4 
    \end{flushleft}
\end{minipage}


\vspace*{0.2cm}
\begin{tabular}{ | M{0.4cm} | M{2cm} | M{6.1cm} | M{0.8cm} | M{0.8cm} | M{0.5cm} | M{6.1cm} |}
    \hline
    \Smiley[1.6] & Время & Ответ и игрок #3 & балл #3 & балл #4 & \tiny{ $\times$ $\leftarrow$ $\rightarrow$} & Ответ и игрок #4 \\
    \hline
    \rotatebox{90}{Задача 1\hspace*{1cm}} & & & & & & \\ [2.1cm]
    \hline
    \rotatebox{90}{Задача 2\hspace*{1cm}} & & & & & & \\ [2.1cm]
    \hline
    \rotatebox{90}{Задача 3\hspace*{1cm}} & & & & & & \\ [2.1cm]
    \hline
    \rotatebox{90}{Задача 4\hspace*{1cm}} & & & & & & \\ [2.1cm]
    \hline
\end{tabular}

\vspace*{0.25cm} 

\begin{minipage}{.55\textwidth}

%Условные значки: \\
$\times$ задача снята \\
$\rightarrow$ переход от #3 к #4 \\
$\leftarrow$ переход от #4 к #3 \\

\begin{tabular}{ | m{5cm} | M{1cm} | M{1cm} |}
\hline
\textbf{} & #3 & #4 \\
\hline
Вольные стрелки & & \\ [0.8cm]
\hline
Итог основного этапа & & \\ [0.8cm]
\hline 
\end{tabular}

\vspace*{1cm}
\textbf{Подписи команд:}\dotfill\\

\vspace*{0.5cm}
\textbf{Подписи судей:}\dotfill\\
\end{minipage}
\hspace*{0.5cm}
\begin{minipage}{.45\textwidth}
\begin{tabular}{ | m{5cm} | M{1cm} | M{1cm} |}
\hline
\textbf{Этап обмена ударами} & #3 & #4 \\
\hline 
Решение задачи соперника & & \\ [0.8cm]
\hline
Штраф за незнание & & \\ [0.8cm]
\hline
Дисциплинарный штраф & & \\ [0.8cm]
\hline
Красота обменной задачи & & \\ [0.8cm]
\hline
\textbf{Итоговый счёт} & & \\ [0.8cm]
\hline
\end{tabular}
\end{minipage}
}



\setlength\parindent{0pt}
\begin{document}


% table 1
\protocol{3}{1.1}{$\rho$}{$\delta$}{Сергей Ламзин}{Леша Смолянинов, Миша Трапезников}

\newpage % table 2
\protocol{3}{1.2}{$\eta$}{$\iota$}{Ваня Адо}{Боря Демешев, Роберт Гринштейн}

\newpage % table 3
\protocol{3}{2.1}{$\beta$}{$\alpha$}{Тая Житина}{Владимир Мезенцев, Вадим Тимергалеев}


\newpage % table 4
\protocol{3}{2.2}{$\varepsilon$}{$\chi$}{Ваня Трофимов}{Егор Адиятов, Андрей Мордасевич}

\newpage % table 5
\protocol{3}{3.1}{$\theta$}{$\kappa$}{Полина Смирнова}{Артем Майдуров, Влада Синицына}

\newpage % table 6
\protocol{3}{3.2}{$\lambda$}{$\nu$}{Марк Прокушкин}{Ян Шапиро}

\newpage % table 7
\protocol{3}{ВК}{$\mu$}{$o$}{Антон Шейкин}{Настя Андреева, Даня Снегур}

\newpage % table 8
\protocol{3}{беседка мед-1}{$\sigma$}{$\tau$}{Лев Каминский}{Максим Федянов}

\newpage % table 9 
\protocol{3}{беседка 1-2}{$\gamma$}{$\xi$}{Вика Луковская?}{Игорь Скрипко}

\newpage % table 10
\protocol{3}{беседка 2-3}{$\psi$}{$\omega$}{Лиза Моисеева}{Михаил Красков}

\newpage % table 11
\protocol{3}{столовая}{$\pi$}{$\varphi$}{Миша Шишацкий}{Ярослав Федулов, Боря Дураков?}


\newpage
{\Huge 
\begin{center}
   \bf ФМТ тур 3\\
    \vspace{3cm}
    \begin{tabular}{| c | c |}
    \hline
    Команды & Аудитория\\ 
    \hline
    {$\rho$}—{$\delta$} & 1.1 \\ % table 1
    \hline 
    {$\eta$}—{$\iota$} & 1.2 \\  % table 2
    \hline
    {$\beta$}—{$\alpha$} & 2.1 \\ % table 3
    \hline
    {$\varepsilon$}—{$\chi$} & 2.2 \\ % table 4
    \hline
    {$\theta$}—{$\kappa$} & 3.1 \\ % table 5
    \hline
    {$\lambda$}—{$\nu$} & 3.2 \\ % table 6
    \hline
    {$\mu$}—{$o$} & ВК \\ % table 7
    \hline
    {$\sigma$}—{$\tau$} & беседка мед-1\\ % table 8
    \hline
    {$\gamma$}—{$\xi$} & беседка 1-2 \\ % table 9
    \hline
    {$\psi$}—{$\omega$} & беседка 2-3 \\ % table 10
    \hline
    {$\pi$}—{$\varphi$} & столовая \\ % table 11
    \hline
    \end{tabular}
    \end{center}
    }
\end{document}


